\documentclass[10pt,a4paper]{article}

\usepackage[T1]{fontenc}
\usepackage[utf8]{inputenc}
\usepackage[ngerman]{babel}
\usepackage{lmodern}
\usepackage[margin=1.55cm,top=1.35cm,bottom=1.35cm]{geometry}
\usepackage[table]{xcolor}
\usepackage{enumitem}
\usepackage{tabularx}
\usepackage{array}
\usepackage{amssymb}
\usepackage{fancyhdr}
\usepackage{microtype}

\definecolor{taskblue}{HTML}{184E77}
\definecolor{tasklight}{HTML}{EAF3F8}
\definecolor{taskgray}{HTML}{F3F4F6}
\definecolor{taskgreen}{HTML}{2D6A4F}

\setlength{\parindent}{0pt}
\setlength{\parskip}{2pt}
\setlist[itemize]{leftmargin=5mm,itemsep=1.2pt,topsep=2pt,parsep=0pt}
\renewcommand{\familydefault}{\sfdefault}
\pagestyle{fancy}
\fancyhf{}
\renewcommand{\headrulewidth}{0pt}
\fancyfoot[L]{\footnotesize SemEval 2027 Task 9 -- Teamarbeitsblatt}
\fancyfoot[R]{\footnotesize Week 4}

\newcommand{\checkbox}{\(\square\)}
\newcommand{\field}[2]{%
  \textbf{#1}\hspace{2mm}\rule{#2}{0.4pt}%
}
\newcommand{\sectiontitle}[1]{%
  \vspace{2mm}%
  {\large\bfseries\color{taskblue} #1}\par
  \vspace{0.5mm}\color{taskblue}\rule{\linewidth}{0.7pt}\color{black}\vspace{1mm}%
}

\begin{document}

\begin{center}
  {\LARGE\bfseries\color{taskblue} SemEval 2027 Task 9}\par
  \vspace{1mm}
  {\Large\bfseries Arbeitsblatt -- Forschungsblock Week 4}\par
  \vspace{1mm}
  {\normalsize 180 Minuten individuelle Arbeit mit gemeinsamer Auswertung}
\end{center}

\vspace{2mm}
\colorbox{tasklight}{%
  \parbox{\dimexpr\linewidth-2\fboxsep\relax}{%
    \field{Name:}{5.2cm}\hfill
    \field{Teamrolle:}{4.3cm}\hfill
    \textbf{Datum:} 03.11.2026
  }%
}

\sectiontitle{1. Grundidee}

\begin{itemize}
  \item Task 9, Subtask 1 ist eine \textbf{Multi-Label-Klassifikationsaufgabe}.
  \item Jeder Absatz kann gleichzeitig mehrere der 14 Framing-Labels besitzen.
  \item Das Modell gibt deshalb einen \textbf{14-dimensionalen binären Vektor} aus.
  \item Beispiel: Der \texttt{Economic}-Klassifikator entscheidet für jeden Absatz:
        \texttt{Economic} -- ja oder nein.
\end{itemize}

\sectiontitle{2. Mögliche Baselines}

\begin{itemize}
  \item Einfache regel- oder häufigkeitsbasierte Baseline.
  \item Ein gemeinsames Modell mit 14 Ausgabewerten und Sigmoid-Aktivierung.
  \item Alternativ 14 unabhängige binäre Klassifikatoren -- einer pro Label.
  \item Das direkte Lernen vollständiger Labelkombinationen ist wahrscheinlich
        ungünstig, weil viele Kombinationen selten sind oder erst im Testset auftreten.
\end{itemize}

\sectiontitle{3. Embedding-basierter Ansatz}

\begin{itemize}
  \item Alle beispielsweise mit \texttt{Political} annotierten Absätze in Embeddings umwandeln.
  \item Embeddings clustern, um typische Untergruppen oder Bedeutungsvarianten zu finden.
  \item Pro Label einen oder mehrere Repräsentanten beziehungsweise Cluster-Zentren bilden.
  \item Neue Texte über ihre Kosinusähnlichkeit zu diesen Label-Repräsentanten klassifizieren.
  \item Ergebnis: eine einfache Nearest-Prototype-Baseline als Vergleich zum trainierten Modell.
\end{itemize}

\sectiontitle{4. Synthetische Trainingsdaten}

\begin{itemize}
  \item Sprachmodelle können zusätzliche Beispieltexte für seltene Labels generieren.
  \item Pro Label gezielt positive und schwierige negative Beispiele erzeugen.
  \item Synthetische Daten nur ergänzend im Training verwenden; echte Daten bleiben der Maßstab.
  \item Evaluation ausschließlich auf echten, unveränderten Daten durchführen.
  \item Generierte Beispiele stichprobenartig auf Labelqualität und Datenlecks prüfen.
\end{itemize}

\sectiontitle{5. Bilder in Text übersetzen}

\begin{itemize}
  \item Ein Vision-Language-Modell beschreibt jedes Nachrichtenbild als Text beziehungsweise Caption.
  \item Optional OCR ergänzen, um sichtbare Schrift, Schilder oder Bildunterschriften zu erfassen.
  \item Bildbeschreibung und Artikelkontext werden anschließend wie normaler Text verarbeitet.
  \item Der Textklassifikator gibt wieder einen 14-dimensionalen Framing-Vektor aus.
\end{itemize}

\vfill
\colorbox{taskgray}{%
  \parbox{\dimexpr\linewidth-2\fboxsep\relax}{%
    \textbf{Leitfrage für heute:} Welche einfache, nachvollziehbare Pipeline liefert uns
    zuerst eine reproduzierbare Vergleichsbasis -- und welche konkrete Verbesserung testen
    wir danach?
  }%
}

\newpage

\begin{center}
  {\LARGE\bfseries\color{taskblue} Hauptaufgaben -- Week 4}\par
  \vspace{1mm}
  {\normalsize Jede Person bearbeitet ein Arbeitspaket und dokumentiert ein überprüfbares Ergebnis.}
\end{center}

\vspace{2mm}
\colorbox{tasklight}{%
  \parbox{\dimexpr\linewidth-2\fboxsep\relax}{%
    \field{Name:}{5.4cm}\hfill
    \field{Arbeitspaket:}{3.8cm}\hfill
    \textbf{Zeit:} 180 Minuten
  }%
}

\sectiontitle{1. Fünf Arbeitspakete -- jede Person übernimmt eines}

\renewcommand{\arraystretch}{1.18}
\begin{tabularx}{\linewidth}{>{\centering\arraybackslash}p{0.8cm} >{\bfseries}p{3.0cm} X >{\raggedright\arraybackslash}p{3.6cm}}
\rowcolor{taskblue}
\color{white} & \color{white}Arbeitspaket & \color{white}Auftrag für heute & \color{white}Ergebnis nach 180 Minuten \\
\checkbox & Literatur und Hypothesen &
3--5 relevante Arbeiten zu Multi-Label-Klassifikation oder Framing Detection finden und deren Methoden vergleichen. &
Quellenliste plus zwei testbare Hypothesen. \\
\rowcolor{taskgray}
\checkbox & Daten und Labels &
Englische Daten prüfen: Labelverteilung, Kein-Frame-Anteil, Multi-Label-Anteil und auffällige Datenprobleme. &
Kurze Statistik plus drei Beobachtungen. \\
\checkbox & Text-Baseline &
14-Output-Modell und 14 binäre Klassifikatoren konzeptionell vergleichen; kleinsten reproduzierbaren Versuch planen oder starten. &
Pipeline, Parameter und erster Laufplan. \\
\rowcolor{taskgray}
\checkbox & Embeddings und Synthese &
Embedding-/Kosinus-Baseline entwerfen; Nutzen von Clustering und synthetischen Beispielen für seltene Labels prüfen. &
Versuchsdesign mit Modellwahl und Risiken. \\
\checkbox & Bild-zu-Text und Evaluation &
Captioning-/OCR-Pipeline für Subtask 2 skizzieren und Varianten für Bild, Caption und Artikelkontext definieren. &
Pipeline-Diagramm plus Vergleichsbaselines. \\
\end{tabularx}

\sectiontitle{2. Zeitplan für die 180 Minuten}

\begin{tabularx}{\linewidth}{>{\bfseries\color{taskblue}}p{2.0cm} X p{5.0cm}}
0--15 Min. & Gemeinsamer Start: Ziel, Datenstand und Arbeitspakete klären. & Paket und Leitfrage eintragen. \\
15--30 Min. & Individuelle Forschungsfrage und erwartetes Ergebnis festlegen. & Hypothese oder Entscheidungskriterium. \\
30--105 Min. & Recherche, Analyse oder erster technischer Versuch. & Quellen, Notizen, Code oder Statistik. \\
105--135 Min. & Ergebnis verdichten und überprüfbar dokumentieren. & Ein teilbares Artefakt vorbereiten. \\
135--165 Min. & Kurzberichte: pro Person maximal sechs Minuten. & Ergebnis, Problem, Empfehlung. \\
165--180 Min. & Gemeinsame Entscheidungen und nächste Schritte. & Zuständigkeit und nächster Termin. \\
\end{tabularx}

\sectiontitle{3. Meine Dokumentation}

\textbf{Meine konkrete Forschungsfrage oder Hypothese:}\par
\vspace{1mm}\rule{\linewidth}{0.35pt}\par\vspace{3mm}

\begin{tabularx}{\linewidth}{X X}
\textbf{Wichtigste Erkenntnisse / Evidenz:} & \textbf{Probleme / offene Fragen:} \\
\vspace{13mm} & \vspace{13mm} \\
\hline
\end{tabularx}

\vspace{2mm}
\textbf{Mein konkretes Ergebnis (Link, Datei, Statistik, Pipeline oder Literaturübersicht):}\par
\vspace{1mm}\rule{\linewidth}{0.35pt}\par\vspace{3mm}

\textbf{Meine Empfehlung an das Team:}\par
\vspace{1mm}\rule{\linewidth}{0.35pt}\par\vspace{3mm}

\colorbox{tasklight}{%
  \parbox{\dimexpr\linewidth-2\fboxsep\relax}{%
    \checkbox\ Ergebnis im Team vorgestellt \hfill
    \checkbox\ Artefakt im Repository/Teamordner abgelegt \hfill
    \checkbox\ Nächster Schritt festgelegt
  }%
}

\end{document}
